%! Choosing and Comparing Models in Context — Unit 2, Lesson 2.7 — Student Packet Cover Sheet
\documentclass[10pt]{article}
\usepackage{atda-article}
\usepackage{atda-boxes}
\usepackage{ltablex}
\keepXColumns

\begin{document}

% Full-bleed royal banner
\begin{tikzpicture}[remember picture, overlay]
  \fill[royal] (current page.north west)
    rectangle ([yshift=-0.9in]current page.north east);
\end{tikzpicture}
\vspace{-0.8in}

\begin{center}
  {\color{white}\textbf{\LARGE Algebra, Trigonometry, and Data Analysis}}\\[4pt]
  {\color{mistmid}\large\textbf{Unit 2: Functions, Transformations, and Graph Analysis}}\\[6pt]
  {\color{white}\textbf{\large Lesson 2.7 \quad Choosing and Comparing Models in Context}}
\end{center}

\vspace{0.22in}
\namedateperiod
\vspace{0.18in}

\begin{learningtargetbox}
\small
\begin{enumerate}[leftmargin=1.5em, itemsep=5pt]
  \item \ldots{} decide whether a \textbf{linear}, \textbf{quadratic}, or \textbf{exponential}
        function best models a table, a graph, or a story --- and say which test decided it.
  \item \ldots{} run the two table tests in order --- \textbf{first differences}, then
        \textbf{second differences}, then \textbf{ratios} --- and know that they only work when the
        inputs go up in \emph{equal} steps.
  \item \ldots{} compare and contrast the three families across every representation I have: parent
        function, table, graph, characteristics, and the story itself.
  \item \ldots{} use a model to answer a question in context, \textbf{graphically} and
        \textbf{algebraically}.
  \item \ldots{} say what a model \emph{cannot} tell me --- where its answers stop being
        trustworthy, and why.
\end{enumerate}
\end{learningtargetbox}

\vspace{0.14in}

\begin{tocbox}
\small
\renewcommand{\arraystretch}{1.5}
\begin{tabularx}{\linewidth}{@{} c l X r @{}}
\rowcolor{mist}
\textbf{\#} & \textbf{Component} & \textbf{Description} & \textbf{Score} \\
\midrule
1 & Warm-Up        & Adding, multiplying, and the three parent shapes  & \blank{1.2cm} \\
2 & Guided Notes   & Two tests run in order, a summary sheet for the whole unit, and a video that
                     cannot really get $3$ billion views                & \blank{1.2cm} \\
3 & Group Activity & Two patches of duckweed, three graphs and three stories, and a coffee cup that
                     hides its own ratio                                & \blank{1.2cm} \\
4 & Exit Ticket    & One gym, two tests, and a building that holds $500$ & \blank{1.2cm} \\
5 & Homework       & Two trucks, two models, and the year they cross    & \blank{1.2cm} \\
\midrule
  & \multicolumn{2}{r}{\textbf{Total}} & \blank{1.2cm} \\
\end{tabularx}
\end{tocbox}

\vspace{0.14in}

\begin{remindbox}
\small
This lesson covers \texttt{AFDA.AF.1c, e, g} and \texttt{AFDA.AF.2h} --- deciding which family
\textbf{models} a situation, comparing the three families across every representation, and using the
model you chose to answer a question. \textbf{One rule runs the whole packet: never name a family
without naming the test.} From a table the test is a row --- first differences, second differences,
or ratios. From a graph it is a feature --- a straight edge, a turning point, a flattening end. And
one warning: \emph{``it curves upward'' rules out linear and nothing else.}
\end{remindbox}

\end{document}
